\documentclass[11pt,a4paper]{article}
\usepackage[T1]{fontenc}
\usepackage{lmodern,amsmath,amssymb,amsthm,microtype}
\usepackage[margin=25mm]{geometry}
\usepackage[colorlinks=true,urlcolor=blue,linkcolor=blue]{hyperref}
\newcommand{\R}{\mathbb R}
\newcommand{\HH}{\mathcal H}
\newtheorem{proposition}{Proposition}
\title{The spectral pentagon interface:\\parameter continuity and regular Weyl representations}
\author{Alan Day \and Chengheng\thanks{AI research assistant (OpenAI).}}
\date{Working note --- 5 November 2026}
\begin{document}\maketitle
\begin{abstract}
We examine two possible explanations of a previously certified nonzero
pentagon residual: an exceptional deformation parameter and a different
regular representation of the same Weyl relations. Dominated convergence
and the standard Stone--von Neumann theorem exclude these explanations
within the specified bounded spectral construction. These are applications
of existing mathematics, conditional on the earlier computer-assisted
bound; they do not independently certify that bound. No quantitative
parameter interval or finite-$G$ physical correlator is obtained.
\end{abstract}

\section{Fixed input and literature interface}
All variables are dimensionless. Write $t>0$ for the deformation parameter
denoted $\hbar$ in Kim--Scarinci (KS) \cite{ks}; it is distinct from the
fixed CCR normalization $\pi$. Let
\[
 \HH=L^2(\R^2,dx\,dy/(2\pi^2)),\quad
 P_j=-\pi i\partial_j,\quad M=e^{iXY/\pi}.
\]
The KS compact-product formula and conjugate pairing give, almost everywhere,
\begin{align}
 F_t(x,y)&=\prod_{j=0}^\infty
 \frac{1+e^{x-(2j+1)\pi t-ity}}{1+e^{x-(2j+1)\pi t+ity}}
 \frac{1+e^{y-(2j+1)\pi/t-ix/t}}{1+e^{y-(2j+1)\pi/t+ix/t}}.\label{prod}
\end{align}
Assign any modulus-one value at a zero denominator. Define by ordinary
joint spectral calculus
\[
 A_t=F_t(X,Y),\quad B_t=F_t(P_y,P_x),\quad
 D_t=A_tB_t-B_tM^{-1}B_tMA_t.
\]
The normalized input is $f(x,y)=\sqrt2e^{-(x^2+y^2)/(2\pi)}$.
Our preceding certificate \cite{cert} gives
\begin{equation}
 \|D_1f\|>b,\qquad b=0.04489.\label{input}
\end{equation}
Equation \eqref{input} is reused, not recomputed here. KS Proposition 6.15
states the pentagon for self-adjoint generators satisfying these Weyl
relations. Its $\Lambda=1$ proof refers back to parameter continuation.
Kim's Proposition 7.4 \cite{kim} explicitly cites the same argument;
it is not an independent check. The following implications concern
exactly \eqref{prod} and these bounded operators.

\newpage
\section{The residual is not confined to the self-dual parameter}
\begin{proposition}
The maps $t\mapsto A_t$, $t\mapsto B_t$, and $t\mapsto D_t$ are strongly
continuous on $(0,\infty)$. Consequently there is an $\varepsilon>0$
such that
\begin{equation}
 |t-1|<\varepsilon\quad\Longrightarrow\quad
 \|D_tf\|>b/2=0.022445.\label{openbound}
\end{equation}
The value of $\varepsilon$ is not estimated here.
\end{proposition}
\begin{proof}
Fix $t_0>0$ and a compact interval $[a,c]\subset(0,\infty)$ around it.
For $|x|,|y|\le R$ the absolute values of the small terms in the two
denominators of \eqref{prod} are bounded by
\[
 e^{R-(2j+1)\pi a},\qquad e^{R-(2j+1)\pi/c}.
\]
Both sequences are summable, uniformly in $t\in[a,c]$. For large $j$
these bounds are below $1/2$; the corresponding ratios differ from
one by at most four times these bounds. Hence the infinite tails
converge uniformly on these compact parameter and spatial sets.

At $t=t_0$, every zero denominator lies in the countable set
\[
 E_{t_0}=\{((2m+1)\pi t_0,(2n+1)\pi/t_0):m,n\in\mathbb Z\}.
\]
This set may overcount cancellations, which is harmless. For a fixed
$(x,y)\notin E_{t_0}$, finitely many initial denominators stay nonzero in
a parameter neighborhood, while the tails already have a uniform
nonzero lower bound. Thus $F_t(x,y)\to F_{t_0}(x,y)$ almost everywhere
as $t\to t_0$. Also $|F_t|=1$ almost everywhere. For every $g\in\HH$,
dominated convergence therefore gives
\[
 \|(A_t-A_{t_0})g\|^2
 =\int |F_t-F_{t_0}|^2|g|^2\,d\mu\longrightarrow0,
 \qquad |F_t-F_{t_0}|^2|g|^2\le4|g|^2.
\]
The fixed Fourier transform diagonalizes $(P_y,P_x)$, so it gives the
same conclusion for $B_t$. Products of strongly continuous uniformly
bounded families are strongly continuous: for example,
\[
 \|(U_tV_t-U_0V_0)g\|
 \le \|U_t\|\|(V_t-V_0)g\|+\|(U_t-U_0)V_0g\|.
\]
Apply this to the two words defining $D_t$; $M$ is fixed and unitary.
Finally use \eqref{input} and continuity of $D_tf$ to obtain
\eqref{openbound}.
\end{proof}
This is strong continuity for a real parameter, not operator-norm
continuity, differentiability, or holomorphic continuation. It includes
both rational and irrational $t$ near one. KS Proposition 6.6(SB1)
distinguishes real Barnes parameter $h$ from $h=it$; its remark on
rational real $h$ cannot be applied to $t=1$ as an exclusion.

\newpage
\section{Changing a regular realization cannot remove the word}
Let $(Q_1,Q_2,R_1,R_2)$ be self-adjoint on a nonzero separable Hilbert
space $\HH_\pi$, with the full Weyl versions of
\[
 [Q_1,R_2]=[Q_2,R_1]=\pi i,
 \quad[Q_1,R_1]=[Q_2,R_2]=[Q_1,Q_2]=[R_1,R_2]=0.
\]
In particular every individual exponential group is strongly continuous.
After the relabeling $q=(Q_1,Q_2)$, $p=(R_2,R_1)/\pi$, these are the
ordinary nondegenerate two-degree-of-freedom CCR. The regularity and
finite-dimensional hypotheses of the Stone--von Neumann theorem are
exactly those in Derezi\'nski \cite{derezinski}, \S3.1--3.2, Theorem 4.

\begin{proposition}
Use the same Borel function $F_t$ and ordinary joint spectral calculus
to define the pentagon word $D_t^\pi$ in this representation. There are
a nonzero multiplicity space $\mathcal K$ and a unitary
$U:\HH\otimes\mathcal K\to\HH_\pi$ such that
\begin{equation}
 U^{-1}D_t^\pi U=D_t\otimes I_{\mathcal K}.\label{unitary}
\end{equation}
Thus $\|D_1^\pi\|>b$, and $\|D_t^\pi\|>b/2$ in the neighborhood above.
\end{proposition}
\begin{proof}
The cited theorem supplies the intertwiner for all Weyl unitaries and
hence their self-adjoint generators and spectral measures. The constant
in the Lebesgue measure of $\HH$ is absorbed in a unitary normalization.
Joint Borel functional calculus commutes with this equivalence.
For the middle factor the closures of the sums satisfy
\[
 (X+P_y,Y+P_x)=M^{-1}(P_y,P_x)M,
\]
as follows first on the Schwartz core by differentiating the chirp,
then by unitary conjugation. In particular the sums strongly commute,
and $F_t(X+P_y,Y+P_x)=M^{-1}B_tM$. This transports every factor of the
pentagon and proves \eqref{unitary}. For any unit vector $k\in\mathcal K$,
\[
 \|D_t^\pi U(f\otimes k)\|=\|D_tf\|,
\]
which proves the claimed lower bounds.
\end{proof}
This directly applies the standard multiplicity theorem; it is not a
new representation-classification result. A simultaneous symplectic
change of the four generators is included. Multiplicity, normalization
of Lebesgue measure, or a mere choice of coordinate polarization does
not change the conclusion. A representation with a different spectrum,
a nonregular Weyl action, or only formal commutators is outside the
stated hypothesis and must be identified explicitly before comparison.

\newpage
\section{What this closes, and what remains open}
\paragraph{Dense-domain changes.}
All factors in the final word are bounded. If the same two bounded
words agree on a norm-dense subspace, continuity extends the equality
to the full Hilbert space. Thus a different dense test domain cannot
repair this particular bounded identity. It can still matter for
intermediate unbounded factors or for a construction defining
different bounded operators.

\paragraph{Null sets and states.}
In the canonical realization each paired joint spectral measure is
unitarily related to Lebesgue measure. Changing $F_t$ only on its
countable exceptional set has no effect, also after \eqref{unitary}.
Some individual states could nevertheless lie in $\ker D_t$; our
conclusion is the failure of an operator identity, not a claim that
every state's residual is nonzero. A proposed constrained physical
sector requires its own observable algebra and inner product.

\paragraph{Evidence boundary and next physical question.}
The positive input \eqref{input} remains computer-assisted. This note
neither independently re-integrates its overlaps nor formally verifies
the numerical library. It derives consequences of that input and
identifies two inadequate explanations of the discrepancy with KS's
literal spectral statement. It does not rule out every quantization of
three-dimensional gravity or any complete de Sitter duality.

The next physical task is to compare an actual observer state's
reflection and gluing prescription with the local algebra used here.
A replacement must specify the changed spectrum, scalar function or
functional calculus, and map its observables and inner product to that
physical problem. It is insufficient to give another formula with the
same name. The finite-$G$ observer and full holographic dictionary
remain open; routine refinement of the previous FFT is not required
by either proposition in this note.

\paragraph{Reading and reproducibility.}
On 5 November 2026 the arXiv records still list KS v2 (21 June 2024)
and Kim v2 (5 July 2023). A targeted search and the checked author and
publisher pages did not locate a relevant correction; this is not an
exhaustive absence claim. The source ledger records exact reading
locations. The previous bound and its file hashes are checked by the
accompanying input audit, which is an integrity check, not another
proof of the overlap enclosures. The present note needs no parameter
scan. Source, PDF, build log, and audit hashes are retained separately.

\begin{thebibliography}{9}\small
\bibitem{ks} H. K. Kim and C. Scarinci,
\emph{A quantization of moduli spaces of 3-dimensional gravity},
\href{https://arxiv.org/abs/2112.13329v2}{arXiv:2112.13329v2},
\S6.3--6.5, Propositions 6.6, 6.11, 6.15.
\bibitem{kim} H. K. Kim, \emph{Three-dimensional quantum gravity from
the quantum pseudo-K\"ahler plane},
\href{https://arxiv.org/abs/2112.13962v2}{arXiv:2112.13962v2},
\S5, Propositions 7.3--7.4.
\bibitem{derezinski} J. Derezi\'nski, \emph{Introduction to
representations of the canonical commutation and anticommutation
relations}, \href{https://arxiv.org/abs/math-ph/0511030v2}{arXiv:math-ph/0511030v2},
\S3.1--3.2, Theorem 4, p.~12.
\bibitem{cert} A. Day and Chengheng, \emph{A posteriori certification
of a bounded pentagon word: Gaussian moments and validated overlap
integrals}, project working note, 1 November 2026;
\texttt{manuscript/residual-certificate/main.tex} and associated evidence.
\end{thebibliography}
\end{document}
